% TODO checklist:
% - [x] Read src/utils/pagination.py for pagination implementation
% - [x] Read src/utils/etag.py for ETag utilities
% - [x] Read src/utils/idempotency.py for idempotency decorator
% - [x] Read src/utils/jsonable.py for JSON serialization utilities
% - [x] Read src/utils/provider_resolver.py for provider resolution
% - [x] Read src/utils/deprecation.py for deprecation framework
% - [x] Referenced Q&A.md for cross-cutting concerns
% - [x] Inspected src/middleware/ for middleware implementations

\section{Platform Engineering Utilities}
\label{sec:platform-utilities}

The Cineca Agentic Platform provides a comprehensive set of reusable utilities that implement common cross-cutting concerns. These utilities ensure consistent behavior across the codebase and reduce boilerplate in endpoint implementations.

\subsection{Pagination and ETags}
\label{subsec:pagination}

The platform implements stateless, token-based pagination with ETag support for efficient caching of paginated responses.

\subsubsection{Stateless Pagination}

Pagination uses offset-based tokens encoded as simple strings, enabling stateless operation without server-side cursor storage:

\begin{lstlisting}[language=Python, caption={Pagination implementation from \texttt{src/utils/pagination.py}.}]
def make_page(
    items: list[Any], 
    page_size: int = 50, 
    page_token: str | None = None
) -> tuple[list[Any], str | None]:
    """Simple stateless pagination over a list.
    
    Args:
        items: Full list of items to paginate
        page_size: Number of items per page (default: 50)
        page_token: Base-10 offset string (e.g., '0', '50')
    
    Returns:
        Tuple of (page_items, next_page_token)
        next_page_token is None when no more pages exist
    """
    try:
        offset = int(page_token) if page_token is not None else 0
    except Exception:
        offset = 0
    
    if page_size <= 0:
        page_size = 50
    
    end = offset + page_size
    page_items = items[offset:end]
    next_token = str(end) if end < len(items) else None
    
    return page_items, next_token
\end{lstlisting}

\subsubsection{Usage in API Endpoints}

\begin{lstlisting}[language=Python, caption={Pagination usage in router endpoints.}]
@router.get("/agent-runs", response_model=PagedResponse[AgentRunSummary])
async def list_runs(
    page_size: int = Query(default=50, ge=1, le=100),
    page_token: str | None = Query(default=None),
    status: str | None = Query(default=None),
    db: Session = Depends(get_db),
) -> PagedResponse[AgentRunSummary]:
    """List agent runs with pagination."""
    
    # Fetch all matching runs (filtered by status if provided)
    runs = repository.list_runs(db, status=status)
    
    # Apply pagination
    page_items, next_token = make_page(runs, page_size, page_token)
    
    return PagedResponse(
        items=page_items,
        next_page_token=next_token,
        total_count=len(runs)
    )
\end{lstlisting}

\subsubsection{ETag Generation}

ETags enable efficient HTTP caching through conditional requests:

\begin{lstlisting}[language=Python, caption={ETag utilities from \texttt{src/utils/etag.py}.}]
import hashlib
import json

def generate_etag(obj: Any, weak: bool = False) -> str:
    """Generate ETag from object's JSON representation.
    
    Args:
        obj: Object to generate ETag for (typically dict/list)
        weak: If True, returns weak ETag (W/"..."), else strong
    
    Returns:
        ETag string in format: "abc123" or W/"abc123"
    """
    json_str = json.dumps(obj, sort_keys=True, default=str, separators=(",", ":"))
    hash_digest = hashlib.sha256(json_str.encode("utf-8")).hexdigest()[:16]
    
    if weak:
        return f'W/"{hash_digest}"'
    return f'"{hash_digest}"'


def validate_etag(if_none_match: str | None, current_etag: str) -> bool:
    """Check if current ETag matches If-None-Match header.
    
    Returns True if ETag matches (caller should return 304).
    Supports weak matching per RFC 7232.
    """
    if not if_none_match:
        return False
    
    if if_none_match.strip() == "*":
        return True
    
    # Parse multiple ETags (comma-separated)
    candidates = [tag.strip() for tag in if_none_match.split(",")]
    
    for candidate in candidates:
        # Strip quotes and W/ prefix for weak comparison
        candidate_val = candidate.lower().replace("w/", "").strip('"')
        current_val = current_etag.lower().replace("w/", "").strip('"')
        
        if candidate_val == current_val:
            return True
    
    return False
\end{lstlisting}

\subsubsection{Context-Aware ETags}

ETags can include route and filter context to ensure proper cache invalidation:

\begin{lstlisting}[language=Python, caption={Context-aware ETag computation.}]
def compute_etag(obj: Any, context: dict[str, Any] | None = None) -> str:
    """Compute weak ETag with optional route/filter context.
    
    Args:
        obj: Primary object to hash (response body)
        context: Route/filter context (ensures ETags vary per route)
    
    Example:
        etag = compute_etag(
            jobs_list,
            context={"route": "user_jobs", "status": "running"}
        )
    """
    hash_input = {"data": obj}
    if context:
        hash_input["context"] = context
    
    raw = json.dumps(hash_input, sort_keys=True, separators=(",", ":"))
    h = hashlib.sha256(raw.encode("utf-8")).hexdigest()
    return f'W/"{h}"'
\end{lstlisting}

\subsection{Idempotency}
\label{subsec:idempotency}

The platform implements idempotency for POST requests to ensure safe retries and prevent duplicate resource creation.

\subsubsection{Idempotency Decorator}

\begin{lstlisting}[language=Python, caption={Idempotency decorator from \texttt{src/utils/idempotency.py}.}]
from functools import wraps
from typing import Any, Callable

def idempotent(
    key_fn: Callable[..., str], 
    ttl: int = 24 * 3600
) -> Callable[[Callable[..., Any]], Callable[..., Any]]:
    """Decorator providing idempotency for FastAPI endpoints.
    
    Args:
        key_fn: Function to generate cache key from request params
        ttl: Time-to-live for cached responses (default: 24 hours)
    
    The decorator:
    1. Extracts Idempotency-Key header from request
    2. If key exists in cache, replays stored response
    3. Otherwise, executes handler and caches response
    """
    def decorator(fn: Callable[..., Any]) -> Callable[..., Any]:
        sig = inspect.signature(fn)
        
        @wraps(fn)
        async def wrapper(*args, **kwargs):
            # Bind parameters for key_fn
            bound = sig.bind_partial(*args, **kwargs)
            
            # Extract Idempotency-Key from headers
            idem_key = _extract_idempotency_key(bound)
            
            if not idem_key:
                # No idempotency requested
                return await fn(*args, **kwargs)
            
            cache_key = key_fn(idem_key, **bound.arguments)
            
            # Check cache
            cached = _get(cache_key)
            if cached is not None:
                # Replay cached response
                return JSONResponse(
                    status_code=cached.get("status", 200),
                    content=cached.get("body"),
                    headers=cached.get("headers", {})
                )
            
            # Execute and cache
            result = await fn(*args, **kwargs)
            envelope = _normalize_response(result)
            _set(cache_key, envelope, ex=ttl)
            
            return result
        
        wrapper.__signature__ = sig
        return wrapper
    
    return decorator
\end{lstlisting}

\subsubsection{Idempotency in Job Creation}

\begin{lstlisting}[language=Python, caption={Idempotency usage for job creation.}]
@router.post("/jobs", response_model=JobCreatedResponse)
@idempotent(
    key_fn=lambda idem_key, **kw: f"jobs:idempotency:{idem_key}",
    ttl=86400  # 24 hours
)
async def create_job(
    request: Request,
    body: JobCreateRequest,
) -> JobCreatedResponse:
    """Create new background job with idempotency support.
    
    Clients should include 'Idempotency-Key' header for safe retries.
    Duplicate requests with the same key return the original response.
    """
    job = await job_service.create(body)
    return JobCreatedResponse(job_id=job.id, status="queued")
\end{lstlisting}

\subsubsection{Two-Tier Idempotency}

For critical operations, the platform implements two-tier idempotency checking:

\begin{lstlisting}[language=Python, caption={Two-tier idempotency check pattern.}]
async def check_idempotency_two_tier(
    key: str,
    redis: Redis,
    db: Session
) -> tuple[bool, dict | None]:
    """Two-tier idempotency: Redis (fast) -> PostgreSQL (authoritative).
    
    Returns:
        (is_duplicate, cached_response)
    """
    # Tier 1: Check Redis cache (fast path)
    cached = await redis.get(f"idempotency:{key}")
    if cached:
        return True, json.loads(cached)
    
    # Tier 2: Check PostgreSQL (authoritative)
    existing = db.query(IdempotencyKey).filter(
        IdempotencyKey.key == key
    ).first()
    
    if existing:
        # Populate Redis cache for future requests
        await redis.set(
            f"idempotency:{key}",
            json.dumps(existing.response),
            ex=86400
        )
        return True, existing.response
    
    return False, None
\end{lstlisting}

\subsection{JSON Utilities}
\label{subsec:json-utils}

The platform provides utilities for safe JSON serialization of complex Python types.

\subsubsection{Type Conversion}

\begin{lstlisting}[language=Python, caption={JSON serialization from \texttt{src/utils/jsonable.py}.}]
from datetime import datetime, date, timezone
from decimal import Decimal
from enum import Enum
from pathlib import Path
from uuid import UUID

def to_jsonable(obj: Any) -> Any:
    """Convert object to JSON-serializable form.
    
    Handles common non-serializable types:
    - datetime/date -> ISO format string (RFC3339 with Z suffix)
    - UUID -> string
    - Decimal -> float
    - Enum -> value
    - set -> list
    - Path -> string
    
    Recursively processes dicts and lists.
    """
    # Primitives pass through
    if obj is None or isinstance(obj, (bool, int, float, str)):
        return obj
    
    # Handle datetime/date
    if isinstance(obj, datetime):
        if obj.tzinfo is None:
            obj = obj.replace(tzinfo=timezone.utc)
        utc_dt = obj.astimezone(timezone.utc)
        return utc_dt.isoformat(timespec='milliseconds').replace('+00:00', 'Z')
    
    if isinstance(obj, date):
        return obj.isoformat()
    
    # Handle UUID
    if isinstance(obj, UUID):
        return str(obj)
    
    # Handle Decimal
    if isinstance(obj, Decimal):
        return float(obj)
    
    # Handle Enum
    if isinstance(obj, Enum):
        return obj.value
    
    # Handle Path
    if isinstance(obj, Path):
        return str(obj)
    
    # Handle set
    if isinstance(obj, set):
        return [to_jsonable(item) for item in obj]
    
    # Handle dict (recursive)
    if isinstance(obj, dict):
        return {key: to_jsonable(value) for key, value in obj.items()}
    
    # Handle list/tuple (recursive)
    if isinstance(obj, (list, tuple)):
        return [to_jsonable(item) for item in obj]
    
    # Handle objects with __dict__
    if hasattr(obj, "__dict__"):
        return to_jsonable(obj.__dict__)
    
    # Fallback: convert to string
    return str(obj)
\end{lstlisting}

\subsubsection{JSONB Column Persistence}

The utility is used when persisting data to PostgreSQL JSONB columns:

\begin{lstlisting}[language=Python, caption={Safe JSONB persistence pattern.}]
def save_agent_run(session: Session, run_data: dict) -> AgentRun:
    """Save agent run with safe JSON serialization."""
    
    # Ensure all nested objects are JSON-serializable
    safe_data = {
        "id": run_data["id"],
        "prompt": run_data["prompt"],
        "output": to_jsonable(run_data.get("output")),
        "steps": to_jsonable(run_data.get("steps", [])),
        "metrics": to_jsonable(run_data.get("metrics", {})),
        "created_at": run_data.get("created_at", datetime.utcnow()),
    }
    
    run = AgentRun(**safe_data)
    session.add(run)
    session.commit()
    return run
\end{lstlisting}

\subsection{Provider Resolution}
\label{subsec:provider-resolution}

The provider resolution utilities handle LLM provider configuration at runtime.

\subsubsection{Base URL Resolution}

\begin{lstlisting}[language=Python, caption={Provider URL resolution from \texttt{src/utils/provider\_resolver.py}.}]
def resolve_provider_base_url(provider: Any) -> str | None:
    """Return effective base URL for provider, applying overrides.
    
    Handles Ollama-specific URL resolution:
    1. Check for explicit OLLAMA_BASE_URL environment override
    2. Fall back to provider config
    3. Auto-detect Docker vs host environment
    """
    base = _coerce_attr(provider, "base_url")
    if not base:
        config = _coerce_attr(provider, "config")
        if isinstance(config, dict):
            base = config.get("base_url")
    
    if is_ollama_provider(provider):
        # Check for environment override
        env_override = getattr(settings, "OLLAMA_BASE_URL", None)
        if env_override:
            base = env_override
        elif not base:
            # Auto-detect appropriate URL
            base = settings.resolve_ollama_base_url()
    
    if base:
        return str(base).rstrip("/")
    return None
\end{lstlisting}

\subsubsection{Ollama Provider Detection}

\begin{lstlisting}[language=Python, caption={Ollama provider detection heuristics.}]
def is_ollama_provider(provider: Any) -> bool:
    """Best-effort detection of Ollama provider."""
    if provider is None:
        return False
    
    candidates = [
        _coerce_attr(provider, "id"),
        _coerce_attr(provider, "name"),
        _coerce_attr(provider, "type"),
        _coerce_attr(provider, "base_url"),
    ]
    
    for value in candidates:
        if value and "ollama" in str(value).lower():
            return True
    
    # Check nested config
    config = _coerce_attr(provider, "config")
    if isinstance(config, dict):
        base = config.get("base_url")
        if base and "ollama" in str(base).lower():
            return True
    
    return False
\end{lstlisting}

\subsubsection{Timeout Configuration}

\begin{lstlisting}[language=Python, caption={Provider-aware timeout configuration.}]
import httpx

DEFAULT_HTTPX_TIMEOUT = httpx.Timeout(
    connect=2.0,
    read=60.0,
    write=10.0,
    pool=5.0,
)

def timeout_for_provider(
    provider: Any, 
    default: httpx.Timeout | None = None
) -> httpx.Timeout:
    """Return appropriate timeout for provider (Ollama-aware).
    
    Ollama requires longer timeouts for model loading.
    """
    if is_ollama_provider(provider):
        secs = int(getattr(settings, "OLLAMA_TIMEOUT_SECS", 60))
        secs = max(secs, 1)
        return httpx.Timeout(
            connect=secs, 
            read=secs, 
            write=secs, 
            pool=5.0
        )
    return default or DEFAULT_HTTPX_TIMEOUT
\end{lstlisting}

\subsubsection{Model ID Mapping}

\begin{lstlisting}[language=Python, caption={Upstream model ID resolution for Ollama.}]
def resolve_upstream_model_id(
    provider: Any,
    resolved_model: str | None,
    requested_model: str | None,
    instance: dict[str, Any] | None,
) -> str | None:
    """Translate logical model IDs to provider-specific IDs.
    
    Ollama uses different model naming conventions, so logical
    IDs like 'phi3-mini-q4' map to Ollama tags like 'phi3:mini'.
    """
    candidate = (instance or {}).get("model_id") or resolved_model
    
    if is_ollama_provider(provider):
        mapping = settings.effective_ollama_model_map
        for key in [candidate, resolved_model, requested_model]:
            if key and key in mapping:
                return mapping[key]
    
    return candidate
\end{lstlisting}

\subsection{Deprecation Framework}
\label{subsec:deprecation}

The platform provides utilities for marking endpoints and features as deprecated.

\subsubsection{Deprecation Headers}

\begin{lstlisting}[language=Python, caption={Deprecation header generation from \texttt{src/utils/deprecation.py}.}]
from datetime import UTC, datetime, timedelta

def deprecation_headers(
    replacement: str | None = None, 
    sunset: str | None = None, 
    *, 
    sunset_days: int = 45
) -> dict[str, str]:
    """Return standardized Deprecation and Sunset headers.
    
    Args:
        replacement: Optional path for Link header (successor-version)
        sunset: Optional explicit Sunset date string
        sunset_days: Days until sunset when sunset not provided
    
    Returns:
        Headers dict including Deprecation, Sunset, and optional Link
    
    Standards:
        - Deprecation header: draft-ietf-httpapi-deprecation-header
        - Sunset header: RFC 8594
        - Link rel="successor-version": RFC 8288
    """
    headers: dict[str, str] = {"Deprecation": "true"}
    
    if sunset:
        headers["Sunset"] = sunset
    else:
        now = datetime.now(UTC)
        sunset_date = now + timedelta(days=sunset_days)
        headers["Sunset"] = sunset_date.strftime("%a, %d %b %Y %H:%M:%S GMT")
    
    if replacement:
        headers["Link"] = f'<{replacement}>; rel="successor-version"'
    
    return headers
\end{lstlisting}

\subsubsection{Usage in Deprecated Endpoints}

\begin{lstlisting}[language=Python, caption={Deprecation in endpoint implementation.}]
from src.utils.deprecation import deprecation_headers

@router.get("/v1/jobs/{job_id}", deprecated=True)
async def get_job_v1(job_id: str) -> Response:
    """Get job details (DEPRECATED).
    
    This endpoint is deprecated. Use /v2/jobs/{job_id} instead.
    """
    # Add deprecation headers to response
    headers = deprecation_headers(
        replacement="/v2/jobs/{job_id}",
        sunset_days=90
    )
    
    job = await job_service.get(job_id)
    
    return JSONResponse(
        content=job.dict(),
        headers=headers
    )
\end{lstlisting}

\subsubsection{OpenAPI Documentation}

Deprecated endpoints are automatically marked in OpenAPI documentation:

\begin{lstlisting}[language=Python, caption={OpenAPI deprecation annotation.}]
@router.get(
    "/v1/legacy-endpoint",
    deprecated=True,  # Shows strike-through in Swagger UI
    tags=["deprecated"],
    summary="Legacy endpoint (use v2 instead)",
    description="""
    **DEPRECATED**: This endpoint will be removed in v3.0.
    
    Use `/v2/new-endpoint` instead.
    
    Sunset date: 2025-06-01
    """
)
async def legacy_endpoint():
    ...
\end{lstlisting}

\subsection{Middleware Components}

The platform includes several middleware components that leverage these utilities.

\subsubsection{Vary Headers Middleware}

\begin{lstlisting}[language=Python, caption={Vary headers for proper caching.}]
def add_vary_headers(response: Response) -> Response:
    """Add Vary headers for proper cache key computation.
    
    Ensures CDN/proxy caches differentiate responses by:
    - Authorization: Different users get different responses
    - X-Tenant-ID: Multi-tenant isolation
    - Accept: Content negotiation
    """
    vary = response.headers.get("Vary", "")
    vary_parts = [v.strip() for v in vary.split(",") if v.strip()]
    
    for header in ["Authorization", "X-Tenant-ID", "Accept"]:
        if header not in vary_parts:
            vary_parts.append(header)
    
    response.headers["Vary"] = ", ".join(vary_parts)
    return response
\end{lstlisting}

\subsubsection{Security Headers Middleware}

\begin{lstlisting}[language=Python, caption={Security headers middleware.}]
class SecurityHeadersMiddleware:
    """Add security headers to all responses."""
    
    async def __call__(self, request: Request, call_next):
        response = await call_next(request)
        
        if settings.ENABLE_SECURITY_HEADERS:
            response.headers["X-Content-Type-Options"] = "nosniff"
            response.headers["X-Frame-Options"] = "DENY"
            response.headers["X-XSS-Protection"] = "1; mode=block"
            
            if settings.ENABLE_HSTS:
                response.headers["Strict-Transport-Security"] = (
                    f"max-age={settings.HSTS_MAX_AGE}; includeSubDomains"
                )
            
            if settings.CSP_POLICY:
                response.headers["Content-Security-Policy"] = settings.CSP_POLICY
        
        return response
\end{lstlisting}

